% engines: lualatex, pdflatex
% expect-warnings: 0
% expect-errors: 0
\documentclass{article}
\usepackage[margin=1cm]{geometry}
\usepackage[spanish]{babel}
\usepackage{../../latex/pgf-go}
\newcommand\case[1]{\par\medskip\noindent\textbf{#1}\par\nopagebreak}
\pagestyle{empty}
\begin{document}

\begin{goban}[9, name = lista]
	\stones[b]{C3, D4, E5}\stones[w]{D3, E4}
	\marks[cross]{F7, G7, F6, G6, H6}
\end{goban}

\case{Una parte: igual que antes}
Punto \pgfgoview{lista}{D4}, rect\'angulo \pgfgoview{lista}{C3:E5}.

\case{Dos partes: solo last (y)}
Las piedras \pgfgoview{lista}{C3, D4} son negras.

\case{Cinco puntos sueltos y una mezcla de punto y rect\'angulo}
Las cruces \pgfgoview{lista}{F7, G7, F6, G6, H6}.
Mezcla: \pgfgoview{lista}{C3, D3:E4, A8 to B9}.

\case{separator y last en la vista y con pgfgoset}
\pgfgoview[separator = {; }, last = { o }]{lista}{C3, D4, E5}.
{\pgfgoset{view = {separator = {\,/\,}, last = {\,/\,}}}\pgfgoview{lista}{C3, D4, E5}.}

\case{Otro idioma: and}
{\pgfgoset{lang = en}\pgfgoview{lista}{C3, D4, E5}.}

\case{Las opciones de goban valen para todas las vistas de la lista}
\pgfgoview[background = false, marks = false]{lista}{F7, G7, F6}.

\end{document}
